% Generated draft: requires mathematical and source review.
\documentclass[11pt]{article}
\usepackage[margin=25mm]{geometry}
\usepackage{fontspec}
\setmainfont{Latin Modern Roman}
\usepackage{amsmath,amssymb,mathtools,unicode-math}
\setmathfont{Latin Modern Math}
\usepackage{enumitem,needspace,xurl,hyperref}
\setlength{\parindent}{0pt}
\setlength{\parskip}{5pt}
\setlength{\emergencystretch}{3em}
\newcommand{\N}{\mathbb{N}}
\newcommand{\Z}{\mathbb{Z}}
\newcommand{\Q}{\mathbb{Q}}
\newcommand{\R}{\mathbb{R}}
\newcommand{\C}{\mathbb{C}}
\newcommand{\F}{\mathbb{F}}
\newcommand{\A}{\mathbb{A}}
\newcommand{\PP}{\mathbb{P}}
\newcommand{\E}{\mathbb{E}}
\newcommand{\eps}{\varepsilon}
\newcommand{\dd}{\,\mathrm{d}}
\begin{document}
\section*{UnsolvedMath ID: 30005224}
\textbf{Source problem code:} catalogue entry
\subsection{Definitions and mathematical statement}
The supplied record asks for optimal lower bounds for two character proportions denoted by \(\theta(G)\) and \(\theta'(G)\), associated with bounds attributed to Thompson and Gallagher. It does not define either proportion or specify the class of groups over which the infima are taken.

A conditional formulation of the requested optimization problem is as follows. Once a nonempty class \(\mathcal G\) of admissible groups and functions \(\theta,\theta'\colon\mathcal G\to[0,1]\) defining the intended proportions have been specified, determine
\[
\alpha=\inf\{\theta(G):G\in\mathcal G\},\qquad
\beta=\inf\{\theta'(G):G\in\mathcal G\}.
\]
Equivalently, determine \(\alpha,\beta\in[0,1]\) such that
\[
\begin{aligned}
&\forall G\in\mathcal G,\quad \theta(G)\geq\alpha,
&&\forall\varepsilon>0\ \exists G\in\mathcal G,\quad \theta(G)<\alpha+\varepsilon,\\
&\forall G\in\mathcal G,\quad \theta'(G)\geq\beta,
&&\forall\varepsilon>0\ \exists G\in\mathcal G,\quad \theta'(G)<\beta+\varepsilon.
\end{aligned}
\]
This is a formal template, not a completed formulation of the source problem: the admissible groups, the characters being counted, the defining conditions, and the denominators of both proportions remain unspecified. In particular, no domain restriction or treatment of degenerate cases can be recovered from the supplied text.

The supplied dated assessment describes universal lower bounds of \(1/3\) for both proportions and evidence toward \(1/2\), while reporting that the optimal bounds remain unproved in general. These are supplied curation claims, not independently verified assertions here. The excerpt refers to Questions 1 and 2 and a main conjecture in a reference labeled [6], but neither that reference's bibliographic entry nor the conjecture's statement is supplied; no conjectural equality is therefore asserted.
\subsection{Short English statement}
Determine the greatest universal lower bounds for the two character proportions attributed to Thompson and Gallagher. The supplied record omits the definitions of both proportions and the admissible class of groups, so only a conditional formulation is possible. Its open status and numerical background are supplied dated assessments requiring review.
\subsection{Sources}
\begin{itemize}
\item Local UnsolvedMath registry, permanent ID \(30005224\), title: \emph{Optimal Lower Bounds for Character Proportions}; exact source problem code: \texttt{catalogue entry}. The record belongs to the Oberwolfach Reports Open Problems collection and records proposed year 2022. No proposer is identified.
\item Supplied report citation: \emph{Character Theory and Categorification} (2023), Oberwolfach Reports, DOI 10.4171/owr/2022/39, pp. 2267--2333. \url{https://doi.org/10.4171/owr/2022/39}. The background describes Oberwolfach Reports as documenting mathematical workshops and their open-problem discussions.
\item Supplied literature source: \emph{Oberwolfach Report 39/2022}. \url{https://ems.press/content/serial-article-files/46974}. The literature assessment was checked on \texttt{2026-08-22}, according to the supplied record. This date and assessment are supplied dated curation, not independently verified source consultation or a current literature determination.
\item The supplied excerpt attributes the lower bounds to Thompson and Gallagher and refers to [6, Questions 1 and 2] and the main conjecture of [6]. The bibliographic identity of [6] and the conjecture's text are absent.
\end{itemize}
\end{document}
